\documentclass[11pt,a4paper]{article}

\usepackage[T1]{fontenc}
\usepackage[margin=24mm]{geometry}
\usepackage{amsmath}
\usepackage{booktabs}
\usepackage{enumitem}
\usepackage[hidelinks]{hyperref}

\newcommand{\code}[1]{\texttt{\detokenize{#1}}}

\title{GNSS\_Processing\_fg Processing Pipeline}
\author{LIGO GNSS Pipeline Summary}
\date{2 August 2026}

\begin{document}
\maketitle

\section{Purpose and scope}

\code{GNSS_Processing_fg.cpp} connects filtered rover observations,
ephemerides, LiDAR/IMU states, receiver-clock states, RINEX base-station
carrier observations, and the GTSAM/iSAM2 graph.  It first initializes the
local-to-ECEF relationship and then incrementally adds measurement and motion
factors, optimizes the float state, optionally fixes carrier ambiguities, and
exports the optimized navigation state.

The primary entry points are \code{processGNSS()}, which validates and buffers
observations, and \code{Evaluate()}, which constructs and optimizes the graph
for an initialized epoch.

\section{Pipeline overview}

\begin{verbatim}
rover observations + ephemerides + LiDAR/IMU state
                         |
                         v
          processGNSSBase(): filter and associate
                         |
              +----------+-----------+
              |                      |
       not initialized          initialized
              |                      |
       fill time window       buffer current epoch
              |                      |
         GNSSLIAlign()             Evaluate()
              |                      |
       coarse ECEF position    initialize graph values
       and local/ECEF align          |
              |                 AddFactor()
              |                      |
              |       +--------------+----------------+
              |       |       |       |       |        |
              |     P/D    clocks   motion  carrier   DD multi-frequency
              |       |       |       |       |        |
              +-------+-------+-------+-------+--------+
                                      |
                                 iSAM2 float update
                                      |
                           ambiguity covariance + LAMBDA
                              |                 |
                            reject            accept
                              |                 |
                       keep float state   add integer priors
                                                |
                                         second iSAM2 update
                                                |
                                      export optimized state
\end{verbatim}

\section{Inputs and graph state}

The processor consumes rover pseudorange, Doppler, carrier phase, uncertainty,
and LLI observations; broadcast ephemerides and ionosphere parameters; the
current LiDAR/IMU state and preintegration result; and, optionally, a local
RINEX 3.x base observation file, ZIP archive, or directory.

\begin{center}
\begin{tabular}{@{}lll@{}}
\toprule
Key & Value & Role \\
\midrule
$A_k$ & position and velocity & Local LiDAR-enabled navigation state \\
$R_k$ & rotation & Platform attitude \\
$O_k$ & motion and biases & Angular rate, acceleration, IMU biases \\
$G_k$ & gravity vector & Per-frame gravity state \\
$F_k$ & ECEF navigation vector & Direct state without LiDAR \\
$B_k$ & four clock biases & GPS/GLO/GAL/BDS receiver clocks \\
$C_k$ & clock drift & Common receiver clock drift \\
$E_0,P_0$ & translation and rotation & Local-to-ECEF alignment \\
$n_j$ & scalar in metres & One reference/subject/band ambiguity arc \\
\bottomrule
\end{tabular}
\end{center}

\section{Observation intake and initialization}

\code{processGNSS()} calls \code{GNSSAssignment::processGNSSBase()} to apply
quality checks, associate ephemerides, and maintain Hatch-smoothed
pseudorange.  Before initialization, epochs with fewer than four valid
satellites are rejected.  Valid GNSS data and simultaneous local pose and
velocity are stored in a \code{WINDOW_SIZE+1} initialization window.

\code{GNSSLIAlign()} requires a complete and temporally continuous window.
Coarse single-point localization estimates ECEF position and constellation
clock bias.  With little local motion, the coarse ECEF point and an ENU
rotation initialize the alignment.  With sufficient motion, coarse ECEF
positions across the window are aligned to the local trajectory.  Failed
localization, excessive time gaps, or failed alignment postpone
initialization rather than seeding an invalid graph.

\code{SetInit()} inserts priors for navigation, alignment, clocks, drift,
biases, and gravity.  An initial iSAM2 update completes initialization.

\section{Per-epoch factor construction}

\code{Evaluate()} computes the epoch interval, selects the LiDAR-enabled or
direct-ECEF representation, initializes the new graph values, and calls
\code{AddFactor()}.

\subsection{Pseudorange, Doppler, and clocks}

For each satellite, transmit time, position, velocity, satellite clock, group
delay, ionosphere, and measurement uncertainty are calculated.  Pseudorange
and Doppler constrain position, velocity, constellation clock bias, and common
clock drift.  Robust noise models are available.  Clock evolution factors
connect consecutive $B_k$ and $C_k$ states.  This legacy path continues to
select the established primary frequency.

\subsection{LiDAR/IMU motion}

The LiDAR-enabled branch normalizes the supplied information matrix and adds a
\code{GnssLioFactor}.  A non-positive diagonal marks the LiDAR contribution
invalid instead of causing division by zero.  The no-LiDAR branch uses IMU
preintegration to connect consecutive direct-ECEF states.

\subsection{Time-differenced rover carrier phase}

The processor pairs the current primary-frequency carrier with an earlier
continuous observation of the same satellite.  The unchanged ambiguity
cancels across time, constraining motion and clock change without an absolute
ambiguity state.  Histories are pruned by graph frame and elapsed time.

\section{Base-station multi-frequency RTK path}

\subsection{RINEX 3.02 ingestion}

\code{BaseStationData} accepts a direct RINEX observation file, a ZIP archive,
or a directory containing multiple files or archives.  It reads ZIP members
without extraction and aggregates consecutive epochs.  The surveyed antenna
coordinate comes from \code{APPROX POSITION XYZ}; multiline
\code{SYS / # / OBS TYPES} records define each constellation's fields.
Carrier values and LLI are decoded for GPS/Galileo L1 and L5, BeiDou B1I and
B2a, and GLONASS L1.  GLONASS channels come from
\code{GLONASS SLOT / FRQ #}.  Special event records are skipped according to
their declared record count.

The time system is normalized to GPST, and duplicate tracking codes at the
same satellite/frequency are collapsed.  Since SSI is signal strength rather
than phase uncertainty, \code{base_carrier_std_cycles} supplies the assumed
standard deviation.  Multiple files must report a consistent ECEF antenna
coordinate.  The configured HKSC archives contain observations from 21 March
2025, irrespective of the enclosing directory's 2026 name.

\subsection{Double-difference construction}

The nearest base epoch must satisfy \code{base_epoch_tolerance}.  Rover and
base records are matched by satellite and physical frequency, then grouped by
constellation and signal band.  A reference satellite is selected separately
for each group.  With subject $s$ and reference $q$, the carrier measurement is

\begin{equation}
  \Delta\nabla L =
  (\lambda_s\Phi_r^s-\lambda_s\Phi_b^s)
  -(\lambda_q\Phi_r^q-\lambda_q\Phi_b^q),
\end{equation}

and the geometric term is

\begin{equation}
  \Delta\nabla\rho =
  (\rho_r^s-\rho_b^s)-(\rho_r^q-\rho_b^q).
\end{equation}

The factor residual is $e_{DD}=\Delta\nabla\rho+n_j-\Delta\nabla L$.
Local-frame and direct-ECEF variants provide analytical Jacobians.

Each ambiguity identity contains its reference satellite, subject satellite,
and band.  L1 and L5 therefore use independent wavelengths, reference
selection, and ambiguity arcs.  Loss of lock, a reference change, or a sample
gap starts a new float ambiguity so an old integer constraint cannot leak
into a new arc.

\section{Float solution and LAMBDA fixing}

The first iSAM2 update produces ambiguities in metres.  Active ambiguities
with sufficient lock duration are converted to cycles together with their
full joint covariance:

\begin{equation}
  \hat a_i=\frac{n_i}{\lambda_i}, \qquad
  Q_{a,ij}=\frac{Q_{n,ij}}{\lambda_i\lambda_j}.
\end{equation}

LAMBDA solves

\begin{equation}
  \breve{\mathbf a}=\arg\min_{\mathbf z\in\mathbb Z^m}
  (\hat{\mathbf a}-\mathbf z)^TQ_a^{-1}
  (\hat{\mathbf a}-\mathbf z).
\end{equation}

Fixing requires the minimum ambiguity count and lock duration, the covariance
precision gate, and the ratio test $s_2/\max(s_1,10^{-12})$ above the configured
threshold.  Partial ambiguity resolution removes the highest-variance state
and retries.  Rejection leaves the float solution unchanged.  Accepted
integers are converted back to metres and inserted as tight fix-and-hold
priors; a second iSAM2 update immediately propagates the constraints into the
navigation state.

\section{State export, fixed-lag management, and reset}

Factor indices are recorded per frame.  Old factors and variables are removed
when the active window exceeds the deletion threshold, and carrier histories
are pruned consistently.  After optional integer fixing,
\code{updateStateFromEstimate()} copies optimized attitude, position,
velocity, biases, and gravity to the external state.

Missing base epochs, unavailable covariance, or failed integer validation
disable only the optional fixed/DD update; the float GNSS/LiDAR/IMU graph can
continue.  \code{Reset()} clears all buffers, histories, ambiguity arcs,
factor identifiers, initialization state, and fix status, and creates a fresh
iSAM2 instance.

\section{Important configuration parameters}

\begin{center}
\begin{tabular}{@{}lll@{}}
\toprule
Parameter & Default & Effect \\
\midrule
\code{use_double_differences} & true & Enable base-station RTK factors \\
\code{use_secondary} & true & Add GPS L2/Galileo E5b/BeiDou B2 \\
\code{use_l5} & true & Add supported L5/E5a/B2a carriers \\
\code{base_station_file} & HKSC directory & RINEX file, ZIP, or directory \\
\code{base_epoch_tolerance} & 0.05 s & Rover/base synchronization gate \\
\code{rtk_ambiguity_gap_tolerance} & 1.5 s & Continuity gate for 1 Hz base epochs \\
\code{base_carrier_std_cycles} & 0.01 cycles & Assumed RINEX phase uncertainty \\
\code{double_difference_sigma} & 0.02 m & DD factor standard deviation \\
\code{ambiguity_prior_sigma} & 100 m & Initial float ambiguity prior \\
\code{enable_integer_fixing} & true & Enable LAMBDA and fix-and-hold \\
\code{lambda_min_ambiguities} & 3 & Dataset partial-fix subset \\
\code{lambda_min_lock_epochs} & 10 & Required continuous lock \\
\code{lambda_ratio_threshold} & 3.0 & Integer validation threshold \\
\code{lambda_max_std_cycles} & 0.25 & Float ambiguity precision gate \\
\code{fixed_ambiguity_sigma_cycles} & 0.001 & Integer-prior sigma \\
\code{rtk_debug} & true & Structured DD/float/LAMBDA logging \\
\code{rtk_debug_epoch_interval} & 10 & Repeated diagnostic period in epochs \\
\bottomrule
\end{tabular}
\end{center}

\section{Interpreting RTK diagnostics (2 August 2026)}

The runtime separates three milestones.  \code{[RTK-DD]} confirms that a base
epoch was synchronized and that DD carrier factors entered the graph; explicit
missing-base and zero-factor counters diagnose failures before optimization.
\code{[RTK-FLOAT]} confirms that ambiguity states exist and reports their
cycle estimates, standard deviations, and lock eligibility.
\code{[RTK-LAMBDA]} reports every validation gate.  Only a line containing
\code{FIX SUCCESS} means an integer RTK solution was accepted and tight
fix-and-hold priors were added.  A \code{NO FIX} or
\code{FLOAT_REJECTED_*} status is still a floating-point solution, with the
specific failed gate shown in the same log stream.

The 2 August replay verified continuous key reuse across 1~Hz base epochs,
correct handling of Galileo LLI bit 2, stable reference selection, and actual
entry into LAMBDA after ten locked epochs.  Three eligible float ambiguities
had standard deviations of roughly 0.85--1.41 cycles, so the conservative
0.25-cycle precision gate correctly retained the float solution.

\section{Current multi-frequency boundary}

Primary, GPS L2/Galileo E5b/BeiDou B2, and L5 are used in the synchronized
base/rover DD path and joint LAMBDA resolution.
The existing pseudorange/Doppler and rover-only time-differenced paths remain
on their primary frequency.  Raw per-frequency integer ambiguities are kept;
no ionosphere-free combination is formed because its combined ambiguity is
not a direct integer in cycles.  GLONASS remains L1-only in this mapping.

\section{Verification status}

The Release executable builds and links successfully.  Eight RTK core tests
verify both DD factor variants and all analytical Jacobians, multi-frequency
selection, LLI masking, RTK arc continuity,
LAMBDA against exhaustive integer search, translation invariance, and invalid
inputs.  Two base-reader tests verify a compact dual-frequency RINEX 3.02
fixture and direct loading of both actual HKSC ZIP archives with
consecutive-hour aggregation.  Further diagnostics print GPS, GLONASS,
Galileo, and BeiDou frequency counts and scan the rover bag.  The synchronized
inputs yield 304 matched epochs, 4449 matched carriers, and 3011 potential DD
factors.  The latest result is fifteen passing LIGO test cases; the catkin
collector reports 30 tests, zero errors, zero failures, and zero skipped tests.

The real rover bag contains no L5 observations.  Its LiDAR topic also uses
\code{livox_ros_driver2/CustomMsg}, while this executable expects
\code{livox_ros_driver/CustomMsg}; therefore a driver adapter remains necessary
before a complete live factor-graph replay can evaluate positioning accuracy
and ambiguity-fix convergence.

\end{document}
